\section{Cumplimiento de los objetivos}
\label{sec:concl_objetivos}
% Extensión estimada: 1.5 páginas, una tabla y tres párrafos (ficha 7.2).
% Fuentes: sec:intro_objetivos (OE1--OE3, tab:intro_objetivos_preguntas), balances de los Caps. 4 y 5,
% sec:disc_limitaciones. El grado de cumplimiento se declara con lo que no se hizo, sin atenuarlo.

El objetivo general era desarrollar y evaluar un sistema, basado en modelos preentrenados y adaptados, capaz de
asistir en el diagnóstico de lesiones de piel y en la generación de sus reportes, y analizar cuán bien clasifica, si
el reporte afirma el diagnóstico correcto y si puede ejecutarse con recursos convencionales. El sistema se construyó
como un \emph{pipeline} de tres componentes, un clasificador, un \gls{knn} sobre la representación de DermLIP y un
generador, y cada una de esas tres preguntas recibió una respuesta medida. La Tabla~\ref{tab:concl_objetivos} resume,
para cada objetivo específico, lo que se hizo, el resultado principal y los capítulos en que se presenta.

\begin{table}[H]
    \centering
    \caption{Objetivos específicos: lo que se hizo, el resultado principal y dónde se presenta.}
    \label{tab:concl_objetivos}
    \small
    \setlength{\tabcolsep}{4pt}
    \begin{tabularx}{\textwidth}{@{}>{\raggedright\arraybackslash}p{2.2cm}
        >{\raggedright\arraybackslash\hspace{0pt}}X >{\raggedright\arraybackslash\hspace{0pt}}X
        >{\raggedright\arraybackslash}p{1.7cm}@{}}
        \toprule
        \textbf{Objetivo} & \textbf{Lo que se hizo} & \textbf{Resultado principal} & \textbf{Capítulos} \\
        \midrule
        OE1. Clasificadores & Cuatro supervisados, dos \glspl{vlm} generativos y DermLIP, evaluados con un mismo
            protocolo y con pruebas estadísticas, los supervisados con tres ejecuciones; umbral de melanoma fijado en validación
            & Swin-Tiny, con F1 macro de $0.733 \pm 0.026$; ningún \gls{vlm} sin una cabeza entrenada alcanza a los
            supervisados; el umbral detecta cuatro de cada cinco melanomas
            & \ref{cp:metodos}, \ref{cp:res_clasif} \\
        \filasep
        OE2. \gls{vlm} para el reporte & Mini-VLM ajustado con LoRA en el codificador visual sobre \emph{SkinCaRe}, comparado con
            LLaVA-Med y SkinVL, sin diagnóstico, en cascada con el clasificador y con el ruteo por concordancia
            & El reporte afirma el diagnóstico que recibe; la consistencia diagnóstica es de 0.849 con la cascada y
            de 0.870 con el ruteo
            & \ref{cp:metodos}, \ref{cp:res_gen} \\
        \filasep
        OE3. Escenarios y costo & Prueba externa en \emph{HIBA}, perturbaciones, calibración, casos dudosos
            separados por el ruteo, métricas del texto y costo de cada componente
            & Pérdida de cerca del 40\,\% de la sensibilidad macro en \emph{HIBA}; riesgo selectivo de 0.085 entre
            los afirmados; Mini-VLM entra en 8.6~GB y escribe un reporte en unos 16 segundos
            & \ref{cp:res_clasif}, \ref{cp:res_gen}, \ref{cp:discusion} \\
        \bottomrule
    \end{tabularx}
\end{table}

El primer objetivo se cumplió en los términos en que se planteó. Los clasificadores de las tres familias se
diseñaron, se implementaron y se evaluaron con la misma partición por lesión y las mismas métricas, y la detección de
melanoma se analizó por separado, con un umbral elegido en validación y no en la prueba. Esa comparación es la que
permite responder la primera pregunta de investigación.

El segundo objetivo se cumplió con una modificación que los resultados justifican. El \gls{vlm} ajustado produce el
reporte e integra la imagen en la descripción, pero no produce por sí mismo un diagnóstico preliminar confiable: sin
un diagnóstico, no afirma de manera confiable la clase de la imagen (Sección~\ref{sec:res_gen_sin_anclaje}). Por eso,
en el sistema final el diagnóstico lo aporta el clasificador y el generador lo transcribe en el reporte. El sistema,
como un todo, produce el diagnóstico preliminar y el reporte, pero el diagnóstico no proviene del \gls{vlm}.

El tercer objetivo se cumplió en parte. Se evaluaron los escenarios que el objetivo nombra: la otra población, con
\emph{HIBA}, los casos en que el sistema no tiene seguridad, con el ruteo, la calidad del texto, con las métricas
automáticas, y el costo computacional, en una GPU de consumo. Quedaron sin evaluar, en cambio, dos aspectos. El
\emph{pipeline} completo no se aplicó a \emph{HIBA}, de modo que la fidelidad del reporte en otra población se infiere
de la caída del clasificador y no se midió. Además, la utilidad práctica se midió solo por el costo: ningún
especialista leyó los reportes, y el efecto del sistema sobre el flujo de trabajo clínico, que el objetivo general
menciona, no se evaluó (Sección~\ref{sec:disc_limitaciones}).
